\documentclass[preview, border={10mm, 5mm, 10mm, 10mm}]{standalone}
\usepackage{tikz, tasks, geometry, xcolor}
\usetikzlibrary{arrows.meta, patterns, calc, intersections, quotes, angles}
\usepackage{amsmath, amssymb, amsfonts, mathtools}
\setlength{\columnsep}{10pt}
\setlength{\columnseprule}{0.4pt}
\usepackage[upright]{fourier}
\usepackage{enumitem}
\geometry{a4paper, margin=1in}
\usepackage{tzplot, pgfplots, kinematikz}
\usepackage{circuitikz}
\ctikzset{resistors/scale=0.75,capacitors/scale=0.75,inductors/scale=0.75}
\usepackage{chemfig}
\usepackage[version=4]{mhchem}
\usepackage[modules=all]{chemmacros}
\usepackage{tkz-euclide}
\usepackage{venndiagram}
\pgfplotsset{compat=1.18}
\everymath{\displaystyle}
\newcommand{\ans}{\textcolor{blue!20!red}{\textit{\quad Ans.}}}
% \renewcommand{\ans}{}  % Uncomment to hide answers
\newenvironment{solution}{\par\noindent\color{red!80!black}$\Rightarrow$\enspace\ignorespaces}{\par}
\newenvironment{alternatesolution}{\par\noindent\color{blue!80!black}$\Rrightarrow$\enspace\ignorespaces}{\par}
\newenvironment{hint}{\par\noindent\color{red!50!black}$\looparrowright$\enspace\ignorespaces}{\par}
\newenvironment{idea}{\par\noindent\color{violet!80!black}$\diamond$\enspace\ignorespaces}{\par}
\newenvironment{remark}{\par\noindent\color{teal!80!black}$\circ$\enspace\ignorespaces}{\par}

% --- Global TikZ style (design uniformity across all diagrams) ---
\tikzset{
    >=latex,
    thick,
    every node/.append style={font=\small},
}
\title{\textsc{}}
\pagestyle{empty}
\begin{document}
\begin{idea}
    \begin{align*}
    \intertext{\textbf{Concept:} Flux quantization for a circular electron orbit in a uniform magnetic field}
    \Phi &= BA \\
    B\pi r^2 &= n\left(\frac{h}{e}\right) \\
    evB &= \frac{mv^2}{r} \\
    \mu &= IA \\
    &= \left(\frac{ev}{2\pi r}\right)\left(\pi r^2\right) \\
    &= \frac{evr}{2} \\
    &= \frac{e^2Br^2}{2m} \\
    &= \frac{nhe}{2\pi m} \\
    \intertext{\textbf{Technique:} Use flux quantization to relate $r$ and $n$, use magnetic force for circular motion to relate $v$ and $r$, then choose the smallest allowed $n$ for the lowest state.}
    \end{align*}
\end{idea}
\hfill $\_\_$
\end{document}
